\documentclass[11pt]{article}
\usepackage[margin=1in]{geometry}
\usepackage{amsmath,amssymb,amsthm}
\usepackage{booktabs}
\usepackage{enumitem}
\usepackage{hyperref}
\hypersetup{colorlinks=true,linkcolor=blue,urlcolor=blue,citecolor=blue}

\theoremstyle{plain}
\newtheorem{theorem}{Theorem}
\newtheorem{proposition}[theorem]{Proposition}
\newtheorem{corollary}[theorem]{Corollary}
\theoremstyle{definition}
\newtheorem{example}[theorem]{Example}
\theoremstyle{remark}
\newtheorem*{remark}{Remark}

\title{2-d Kernels in Graphs and Digraphs: Every Result, Proved or Disproved}
\author{Reference compendium}
\date{August 21, 2026}

\begin{document}
\maketitle
\vspace{-1.5em}

Throughout, a \textbf{2-d kernel} of a digraph $D$ is a set $S \subseteq V$ that is
\emph{independent} (no arc, in either direction, joins two vertices of $S$) and
\emph{2-dominating} (every $v \notin S$ has arcs to at least 2 distinct vertices of
$S$). Undirected graphs are treated as symmetric digraphs (each edge becomes arcs in
both directions). $N^+(v)$ is the out-neighbourhood, $\delta^+(D)$ the smallest
out-degree in $D$. Deciding 2-d kernel existence on general graphs is NP-complete
(J.~A.~Telle, \emph{Nordic J. Comput.} \textbf{1} (1994) 157--171) — every
polynomial-time result below is therefore a genuine dividing line, not a formality.

\section*{Overview of all results}

\begin{enumerate}[leftmargin=2.2em]
\item \textbf{Bipartite digraphs (Theorem A)} — trivial 2-d kernel whenever one side
has out-degree $\ge 2$. \textit{Proved.}
\item \textbf{Strongly connected, all-even-cycle digraphs (Theorem B)} — two disjoint
2-d kernels guaranteed. \textit{Proved.}
\item \textbf{Dropping strong connectivity} — the hypothesis is not removable.
\textit{Disproved} (explicit 6-vertex counterexample).
\item \textbf{Chordal graphs} — 2-d kernel existence and uniqueness decidable in
polynomial time. \textit{Proved.}
\item \textbf{Simplicial graphs} — decided by a simpler argument on an overlapping
but different class. \textit{Proved.}
\item \textbf{Split graphs} — full existence characterization, two cases.
\textit{Proved.}
\item \textbf{Bipartite graphs} — conjecture that the same shortcut always decides
them. \textit{Disproved} (8-vertex counterexample).
\item \textbf{$k$-degenerate graphs} — attempted generalization of the chordal
result. \textit{Disproved} ($C_5$ counterexample at degeneracy 2).
\item \textbf{Chordal-underlying digraphs} — the chordal result lifts to digraphs
with no extra restriction. \textit{Proved.}
\item \textbf{Oriented graphs, tight bound} — no 2-d kernel below 8 vertices once
digons are forbidden. \textit{Proved.}
\item \textbf{Oriented graphs, extremal count} — exactly 2 isomorphism classes at
the minimum, under a necessary added hypothesis. \textit{Proved (hypothesis-qualified)}.
\item \textbf{DAGs} — unique 2-d kernel, decidable in linear time. \textit{Proved.}
\item \textbf{Subdivisions} — a free infinite family that always has one.
\textit{Proved.}
\item \textbf{Cubic graphs and girth} — whether avoiding short cycles helps.
\textit{Null result} — no supporting trend found.
\end{enumerate}

\section{Bipartite digraphs}

\begin{theorem}[A]
If the underlying graph of a digraph $D$ is bipartite with parts $P, Q$, and every
vertex of $Q$ has $d^+ \ge 2$, then $P$ is a 2-d kernel.
\end{theorem}

\begin{proof}
Bipartiteness means every arc runs between $P$ and $Q$, so $P$ is automatically
independent (no arc stays inside $P$). Every $v \in Q$ has all its out-arcs landing
somewhere, and none of them can land back inside $Q$, so all of $v$'s out-arcs land
in $P$ — at least 2 of them.
\end{proof}

\section{Strongly connected digraphs where every cycle is even}

\begin{theorem}[B]
If $D$ is strongly connected, every directed cycle of $D$ has even length, and
$\delta^+(D) \ge 2$, then $D$ has two disjoint 2-d kernels.
\end{theorem}

\begin{proof}
\emph{A closed directed walk decomposes into cycles.} If $W = w_0,\dots,w_L$
($w_0=w_L$) is a closed directed walk, either $w_0,\dots,w_{L-1}$ are all distinct (so
$W$ is itself a simple cycle of length $L$), or some $w_a = w_b$ with $a<b\le L-1$,
splitting $W$ into two strictly shorter closed walks — of lengths $b-a$ and
$L-(b-a)$ — whose cycle decompositions (by induction) combine to one for $W$, with
lengths still summing to $L$.

Fix any vertex $r$. \emph{Claim: every directed path from $r$ to a vertex $v$ has the
same length parity.} If $P, Q$ were two $r\to v$ paths of different parity, take any
directed path $R$ from $v$ back to $r$ (strong connectivity). $P{+}R$ and $Q{+}R$ are
closed walks of lengths $|P|{+}|R|$, $|Q|{+}|R|$; by the decomposition fact, both
lengths are sums of even cycle-lengths, hence even — forcing $|P|\equiv|Q|\pmod 2$,
a contradiction. So every vertex $v$ has a well-defined parity $\mathrm{par}(v)$.

Let $V_0, V_1$ be the two parity classes (partitioning $V$, since $D$ is strongly
connected). Every arc $u\to v$ flips parity — extend an $r\to u$ path by this arc to
get an $r\to v$ path one longer — so no arc joins two same-parity vertices ($V_0,
V_1$ are independent), and every out-neighbour of a $V_i$-vertex lies in $V_{1-i}$.
So $S=V_{1-i}$ satisfies: every vertex outside $S$ (i.e.\ of $V_i$) has \emph{all} of
its $\ge\delta^+(D)\ge2$ out-neighbours inside $S$ — 2-domination, automatically. Both
$V_0$ and $V_1$ are 2-d kernels, and they're disjoint since they partition $V$.
\end{proof}

\section{Strong connectivity cannot be dropped}

\begin{example}
Vertices $\{1,2,3,4,a,b\}$. Symmetric 4-cycle on $1,2,3,4$ (arcs both ways between
each consecutive pair), plus $a\to1$, $a\to2$, $a\to b$, $b\to2$, $b\to3$. Every
directed cycle is even (the digons are length-2, the outer ring is length 4, and $a,b$
have nothing pointing into them so can't sit on any cycle). Every vertex has
out-degree $\ge2$. But nothing reaches $a$ or $b$ — not strongly connected — and this
digraph has \textbf{no} 2-d kernel.
\end{example}

\begin{proof}[Why no 2-d kernel exists]
Within $\{1,2,3,4\}$ alone (an even cycle) the only two candidates are $\{1,3\}$ and
$\{2,4\}$. Try $S \supseteq \{1,3\}$: $a$ is then forced \texttt{OUT} (it has an arc
to $1\in S$, and joining $S$ itself would violate independence against that arc), so
$a$ needs 2 out-neighbours in $S$ from $\{1,2,b\}$; it has $1$, and needs $2$ or $b$
— but $2\notin S$ here, so $b$ must join $S$. But $b\to3$ is an arc and $3\in S$, so
$b$ joining $S$ violates independence. Contradiction. The case $S\supseteq\{2,4\}$
fails symmetrically ($a$ again forced \texttt{OUT}, needing $b\in S$, but $b\to2$ and
$2\in S$ blocks it). Both possibilities fail.
\end{proof}

\emph{Checked exhaustively}: among all 126 digraphs on $\le6$ vertices with
$\delta^+\ge2$ that are all-cycles-even but \textbf{not} strongly connected, 84
(66.7\%) still lack a 2-d kernel or have one only by luck — this is a typical failure
mode, not a contrived exception. Among the 81 that \emph{are} strongly connected and
all-cycles-even, all 81 have one, exactly as Theorem B guarantees.

\section{Chordal graphs}

A vertex $x$ is \textbf{simplicial} if its neighbourhood $N(x)$ is a clique. Every
non-empty chordal graph has one, and every induced subgraph of a chordal graph is
chordal — both classical facts, used but not re-derived here.

\begin{proposition}[Soundness of the two rules used below]
Run two rules to a fixed point, starting all vertices \texttt{UNKNOWN}: \textbf{(Spread)}
if $v$ is \texttt{IN}, mark every neighbour \texttt{OUT}; \textbf{(Force)} if $v$ is
\texttt{UNKNOWN} and its non-\texttt{OUT} neighbours contain no independent pair
(i.e.\ form a clique, possibly of size $\le1$), mark $v$ \texttt{IN}. Then every
status this fixes holds in \emph{every} 2-d kernel of $G$, not just some candidate.
\end{proposition}

\begin{proof}
Induction on the order vertices get fixed. If $v$ is fixed \texttt{OUT} via Spread,
it was triggered by a neighbour $u$ already \texttt{IN}; by induction $u\in S$ for
every 2-d kernel $S$, so $v\notin S$ (independence). If $v$ is fixed \texttt{IN} via
Force: take any 2-d kernel $S$ and suppose $v\notin S$. Every neighbour of $v$ already
\texttt{OUT} is (by induction) excluded from $S$, so $S\cap N(v)$ sits entirely
inside $v$'s current non-\texttt{OUT} neighbours. Since $v\notin S$, 2-domination
needs $|S\cap N(v)|\ge2$ — but $S$ is independent, so those $\ge2$ vertices are a
non-adjacent pair inside $v$'s non-\texttt{OUT} neighbourhood, contradicting the very
condition that made Force fire. So $v\in S$.
\end{proof}

\begin{theorem}
On a chordal graph, running (Spread) and (Force) to a fixed point never leaves a
vertex \texttt{UNKNOWN}. Consequently a chordal graph has at most one 2-d kernel, and
2-D KERNEL is solvable in polynomial time on chordal graphs.
\end{theorem}

\begin{proof}
Suppose some vertex remains \texttt{UNKNOWN}; let $W\ne\emptyset$ be all such
vertices. No $v\in W$ has an \texttt{IN} neighbour (Spread would already have fired),
so $v$'s non-\texttt{OUT} neighbours are exactly $N(v)\cap W$. $G[W]$ is an induced
subgraph of a chordal graph, hence chordal, hence has a simplicial vertex $x$: $N(x)
\cap W$ is a clique — exactly $x$'s current non-\texttt{OUT} neighbourhood — so Force
fires on $x$, contradicting the fixed point. Hence $W=\emptyset$. By the Proposition,
the resulting set is the \emph{only} possible 2-d kernel of $G$.
\end{proof}

\emph{Confirmed exhaustively} on all 17{,}175 chordal graphs up to 9 vertices: 0
\texttt{UNDECIDED}, maximum 2-d kernel count 1.

\emph{Literature note}: a closely related dichotomy (Golovach--Kratochv\'il, WG 2007)
covers this shape of problem for finite domination thresholds; our threshold is
cofinite ($\{2,3,4,\dots\}$), so this result is presented independently, not as a
confirmed corollary of theirs — their exact scope could not be verified past the
abstract.

\section{Simplicial graphs}

A graph is a \textbf{simplicial graph} if every vertex lies in $N[x]=N(x)\cup\{x\}$
for some simplicial $x$ — weaker than chordal, and a different (overlapping) class.

\begin{theorem}
On a simplicial graph, rules (Spread) and (Force) — in fact just the simpler pair,
mark every simplicial vertex \texttt{IN}, then spread — decide the whole graph, and
there is at most one 2-d kernel.
\end{theorem}

\begin{proof}
Every simplicial vertex has a clique neighbourhood, so Force marks all of them
\texttt{IN} immediately (before anything else needs to happen). Every other vertex,
by definition of the class, is adjacent to some simplicial vertex — now \texttt{IN}
— so Spread marks it \texttt{OUT}. Nothing is left \texttt{UNKNOWN}.
\end{proof}

\emph{Confirmed exhaustively} on 14{,}131 simplicial graphs: 0 disagreements, max
count 1. (Attach a pendant to every vertex of a 4-cycle for an example that is
simplicial but \emph{not} chordal — the 4-cycle itself has no chord.)

\section{Split graphs}

A \textbf{split graph} has vertices partitioned into a clique $C$ ($|C|\ge2$) and an
independent set $I$, with arbitrary edges between them.

\begin{theorem}
A split graph $G=(C,I)$ has a 2-d kernel if and only if either
\textup{(a)} every vertex of $C$ has $\ge2$ neighbours in $I$, or
\textup{(b)} exactly one $c_0\in C$ has \emph{no} neighbour in $I$, and every other
vertex of $C$ has $\ge1$ neighbour in $I$.
\end{theorem}

\begin{proof}
Since $C$ is a clique, any independent $S$ has $|S\cap C|\le1$. Two cases, and they
are exhaustive.

\emph{$S\cap C=\emptyset$, so $S\subseteq I$.} Since $I$ is internally edge-free, any
$v\in I\setminus S$ is adjacent to nothing in $S$ at all — so $S$ isn't maximal
unless $S=I$ exactly. (Every 2-d kernel is a maximal independent set, since anything
outside it needs arcs into it.) So $S=I$ is forced. Domination then requires every
$c\in C$ (all of which sit outside $S=I$) to have $\ge2$ neighbours in $I$ — case (a).

\emph{$S\cap C=\{c_0\}$ for some specific $c_0$.} Take any $v\in I$ adjacent to
$c_0$: $v$'s only possible neighbours are in $C$ (since $I$ is edge-free), and the
only $C$-vertex in $S$ is $c_0$ itself — so $v$ has at most 1 neighbour in $S$,
never enough for 2-domination. So no such $v$ can exist: $c_0$ must have \emph{zero}
neighbours in $I$. With that settled, every $v\in I$ is non-adjacent to $c_0$, so (by
the maximality argument above, now applied with $c_0$ already in $S$) all of $I$
must join $S$ too: $S=\{c_0\}\cup I$. Domination of the remaining vertices, each
$c'\in C\setminus\{c_0\}$: $c'$ is adjacent to $c_0$ (clique) — 1 neighbour in $S$
already — and needs 1 more from $S\cap I=I$, i.e.\ $c'$ needs $\ge1$ neighbour in
$I$. That is exactly case (b).
\end{proof}

\emph{Confirmed}: zero mismatches against brute force over 20{,}000 random split
graphs up to 9 vertices.

\section{Bipartite graphs — a refuted conjecture}

At $n\le7$, all 150 bipartite graphs in the census were fully decided by (Spread) and
(Force) — as clean a pattern as chordal's. It breaks at $n=8$.

\begin{example}
The 8-vertex graph with canonical encoding \texttt{G?KsZ\_}, parts $\{0,4,5,7\}$ and
$\{1,2,3,6\}$, has two pendant (degree-1) vertices $0$ and $1$ whose unique
neighbours ($6$ and $7$) lie in \emph{different} bipartition classes. It has no
2-d kernel at all.
\end{example}

\begin{proof}[Why forcing cannot resolve it]
Both pendants fire Force immediately (trivially, a degree-1 vertex's single neighbour
can't contain an independent pair), forcing $6,7$ \texttt{OUT} via Spread. Now
$N(6)\setminus\text{OUT}=\{0,4,5\}$ and $N(7)\setminus\text{OUT}=\{1,2,3\}$ — each a
3-element set containing a non-adjacent pair, so Force cannot fire on either. Nothing
else changes; several vertices remain \texttt{UNKNOWN} forever, even though (checked
by brute force) this graph in fact has no 2-d kernel. This does not contradict Wloch's
even-distance theorem for bipartite graphs (Australas. J. Combin.\ \textbf{53}
(2012), Thm.\ 2.4) — it is exactly the situation that theorem's hypothesis excludes.
\end{proof}

This is the project's clearest lesson: two conjectures indistinguishable on evidence
up to 7 vertices turned out to be one theorem and one false statement. Only pushing
the exhaustive range further — not more theorizing — told them apart.

\section{$k$-degenerate graphs — a refuted generalization}

A graph is \textbf{$k$-degenerate} if it can be reduced to nothing by repeatedly
removing a vertex of degree $\le k$. Every chordal graph is, in particular, built
this way at low degeneracy — so the natural next question was whether the chordal
argument survives weakening ``chordal'' to ``low degeneracy.''

\begin{example}
$C_5$, the 5-cycle: every vertex has degree exactly 2, so degeneracy 2. But no
vertex of $C_5$ is simplicial — there isn't a single triangle in it, so no vertex's
neighbourhood is a clique.
\end{example}

\begin{proof}[Why the chordal argument breaks here]
The chordal proof's only lever is: the undecided region, at a fixed point, is always
chordal, hence always has a simplicial vertex to force in. Degeneracy 2 alone
guarantees a low-degree vertex at every stage of removal, but says nothing about
that vertex's neighbourhood being a clique. $C_5$ realizes exactly this gap: Force
never fires on any vertex (every vertex's two neighbours are non-adjacent to each
other), so the whole graph stays \texttt{UNDECIDED} — yet $C_5$, an odd cycle, in
fact has no 2-d kernel at all (Section 12's cycle fact).
\end{proof}

\emph{Confirmed exhaustively}: at degeneracy $\le1$ (forests — already subsumed by
the chordal theorem, since forests are chordal), forcing is complete, 0
disagreements. At degeneracy 2, checked across 35{,}739 graphs: 3{,}388 disagreements
(9.5\%), maximum 2-d kernel count rises to 4. The conjecture fails at the very next
step past forests.

\section{Chordal-underlying digraphs}

\begin{theorem}
If the underlying graph of a digraph $D$ is chordal, then (Spread) and (Force) —
using $N^+(v)$ in place of $N(v)$ — decide $D$ completely, so $D$ has at most one
2-d kernel, with no restriction on out-degree at all.
\end{theorem}

\begin{proof}
The one place this could plausibly break is that Force, for a digraph, tests
$N^+(v)$ (out-neighbours only) rather than the full $N(v)$ that simpliciality is
defined on. But $N^+(v)\subseteq N(v)$ always — the out-neighbourhood is a subset of
the full neighbourhood — and any subset of a clique is still a clique. So if $v$ is
simplicial ($N(v)$ a clique), $N^+(v)$ is automatically a clique too, regardless of
which arcs happen to point outward. With that one step secured, the chordal-graph
proof (Section 4) carries over completely unchanged: the undecided region is chordal,
has a simplicial vertex, and its out-neighbourhood within that region is a subset of
a clique, hence a clique, firing Force.
\end{proof}

\emph{Confirmed exhaustively} on 251{,}991 digraphs (underlying graph chordal, $\le6$
vertices, $\delta^+\ge2$): 0 disagreements, plus 20{,}000 further random digraphs with
\emph{no} out-degree restriction at all, confirming the stronger unrestricted
statement above.

\section{Oriented graphs: the tight bound}

An \textbf{oriented graph} has no \textbf{digons} (no pair of opposite arcs between
the same two vertices) — the opposite extreme from treating a plain edge as two
arcs.

\begin{theorem}
Let $D$ be oriented with $\delta^+(D)\ge2$. If $D$ has a 2-d kernel, $D$ has at least
8 vertices, and 8 is attained.
\end{theorem}

\begin{proof}
Let $J$ be a 2-d kernel, $j=|J|$, $B=V\setminus J$, $t=|B|>0$. Writing $d_x$ for the
number of $B$-vertices pointing into $x\in J$: 2-domination gives
$\sum_{x\in J}d_x\ge2t$. Fix $x\in J$: since $J$ is independent, all of $x$'s
out-arcs go to $B$; since there are no digons, $x$ cannot point at any $B$-vertex
that already points at it, so $x$'s out-arcs are confined to the $t-d_x$ vertices of
$B$ that don't — and since $d^+(x)\ge2$, this forces $d_x\le t-2$. Combining,
$2t\le\sum d_x\le j(t-2)$, so $j\ge\frac{2t}{t-2}$ (requiring $t\ge3$). Minimising
$n=j+t$: $t=3\Rightarrow n\ge9$; $t=4\Rightarrow n\ge8$ (minimum); $t=5,6\Rightarrow
n\ge9$; $t\ge7\Rightarrow n\ge10$ and rising. So $n\ge8$ always, tight only at
$t=j=4$.
\end{proof}

\begin{example}
$J=\{0,1,2,3\}$, $B=\{4,5,6,7\}$, arcs $(0,5)(0,6)(1,6)(1,7)(2,4)(2,7)(3,4)(3,5)$ and
$(4,0)(4,1)(5,1)(5,2)(6,2)(6,3)(7,0)(7,3)$: every vertex out-degree exactly 2, no
digon, and both $J$ and $B$ are 2-d kernels.
\end{example}

\section{Oriented graphs: the extremal count}

\begin{proposition}
Among 8-vertex oriented graphs in which \emph{every} vertex has out-degree
\emph{exactly} 2 (not merely $\ge2$) and which have a 2-d kernel of size 4, there are
exactly 2 isomorphism classes (sizes 72 and 18 of 90 labelled instances), both
balanced orientations of $K_{4,4}$.
\end{proposition}

\begin{proof}
At $t=j=4$, both inequalities above collapse to equalities: $\sum d_x=8$ exactly,
and since each $d_x\le2$ with 4 terms summing to 8, every $d_x=2$ — forcing
$d^+(x)=2$ exactly for every $x\in J$ (this much is forced by minimality alone, no
extra hypothesis). Likewise every $b\in B$ sends exactly 2 arcs into $J$ (same
averaging argument). What is \emph{not} forced is that $b$ has no further out-arcs
elsewhere — $\delta^+\ge2$ is a lower bound only. Under the added hypothesis that
every vertex's out-degree is exactly 2, a $B$-vertex has no budget left once its 2
arcs to $J$ are placed, so $B$ has no internal arcs, and the whole graph reduces to
one choice per $B$-vertex (which 2 of the 4 $J$-vertices it targets, subject to
every $J$-vertex being targeted exactly twice; the $J\to B$ arcs are then forced).
Exhaustive enumeration of the $6^4=1296$ labelled choices keeps exactly 90 balanced
ones; grouping by isomorphism gives exactly 2 classes, sizes 72 and 18 — verified
directly by computer search.
\end{proof}

\begin{remark}
This added hypothesis is necessary, not cosmetic: take the Example above and add one
further arc entirely within $B$, say $4\to5$ (creates no digon, since $5\to4$ was
never present). This is still oriented, 8 vertices, $\delta^+\ge2$, and $J$ remains
a valid 2-d kernel (its independence and domination depend only on arcs touching $J$).
This graph has 17 arcs, so it cannot be isomorphic to any of the 90 sixteen-arc
graphs above — a third, uncatalogued extremal example. So the Proposition's count of
2 is exactly as qualified, not a count of every 8-vertex example satisfying only
Theorem 10's own hypotheses.
\end{remark}

\section{DAGs}

\begin{theorem}
A DAG has at most one 2-d kernel, decidable in linear time.
\end{theorem}

\begin{proof}
Fix a topological order $v_1,\dots,v_n$ (arcs run low-to-high index) and process in
reverse. When $v_i$'s turn comes, every out-neighbour has already been permanently
fixed \texttt{IN}/\texttt{OUT}. Let $k$ be the count of $v_i$'s out-neighbours
already \texttt{IN}.
\begin{itemize}[nosep]
\item $k=0$: \texttt{OUT} is impossible (needs $k\ge2$); \texttt{IN} is consistent
(no \texttt{IN} out-neighbour to conflict with). Set $v_i:=\texttt{IN}$.
\item $k=1$, via out-neighbour $u$: \texttt{IN} is impossible (arc $v_i\to u$ with
$u$ already \texttt{IN} would put two adjacent vertices both in the set);
\texttt{OUT} is impossible ($k<2$). \emph{Neither status is available} — no
2-d kernel of $D$ exists. Halt.
\item $k\ge2$: \texttt{IN} is impossible (already has an \texttt{IN}
out-neighbour); \texttt{OUT} is consistent ($k\ge2$ satisfies domination). Set
$v_i:=\texttt{OUT}$.
\end{itemize}
The base case $v_n$ is a sink ($N^+(v_n)=\emptyset$), so $k=0$ trivially:
$v_n:=\texttt{IN}$. If the sweep completes, let $S=\{v:\texttt{IN}\}$. Independence
holds \emph{by construction}: any arc $v_i\to v_j$ with $v_i$ set \texttt{IN} was set
via the $k=0$ case, meaning $v_j$ was not \texttt{IN} at that time — so $v_j$ is
\texttt{OUT}, not in $S$. Domination of every \texttt{OUT} vertex holds directly from
the $k\ge2$ condition that placed it there. So $S$ is a genuine 2-d kernel, and since
every step is forced with no branching, the only one.
\end{proof}

\emph{Confirmed} on all 20{,}237 DAGs up to 7 vertices: max count 1, 0
\texttt{UNDECIDED}. The fraction of DAGs \emph{with} a 2-d kernel falls steadily
(100\%, 75\%, 50\%, 29.5\%, 16.9\% at $n=1,4,5,6,7$) — this is the growing chance
that some vertex lands at $k=1$ somewhere in the sweep, a specific structural event,
not a vague ``check that might fail.''

\section{Subdivisions}

The \textbf{subdivision} $S(G)$ replaces every edge $\{u,v\}$ with a fresh vertex $w$
and edges $\{u,w\},\{w,v\}$.

\begin{theorem}
For every graph $G$, $V(G)$ is a 2-d kernel of $S(G)$.
\end{theorem}

\begin{proof}
No two original vertices are adjacent in $S(G)$ (every original edge is now a path
through a subdivision vertex), so $V(G)$ is independent. Every subdivision vertex's
neighbourhood is exactly the two endpoints of the edge it replaced — both in $V(G)$
— exactly 2, satisfying 2-domination.
\end{proof}

A free infinite family, checked directly on all 1253 atlas graphs (largest
subdivision checked: 28 vertices).

\section{Cubic graphs and girth — a null result}

Does forbidding short cycles make a 2-d kernel more likely, among 3-regular graphs?
Seeded random samples, $n=10,\dots,20$:

\begin{center}
\begin{tabular}{@{}lrrr@{}}
\toprule
family & classes & with a 2-d kernel & fraction \\
\midrule
cubic, unrestricted & 293 & 160 & 0.546 \\
cubic, triangle-free & 240 & 131 & 0.546 \\
cubic, girth $\ge5$ & 158 & 79 & 0.500 \\
\bottomrule
\end{tabular}
\end{center}

No trend: the fraction sits near one half regardless of girth. Not a theorem, not a
counterexample — just no evidence found for the idea, at least at degree 3.

\end{document}
